\section{Conformal Equivalence to Circular Geometry}
\label{sec:conformal-equivalence}

The closed-form solution of Sec.~\ref{sec:exact-solution} admits a simple
geometrical interpretation.  Optical metrics and differential-geometric
methods have also been used to formulate atmospheric refraction
\cite{kropla1992differential}.  For the Earth-scale choice $\kappa=1/R$, the
optical
line element associated with the exponential profile is
\begin{equation}
    d\ell_{\mathrm{opt}}^2
    =
    n_0^2 e^{2h/R}
    \left(dx^2+dh^2\right).
    \label{eq:exponential-optical-metric}
\end{equation}
Introduce the coordinates
\begin{equation}
    r=R e^{h/R},
    \qquad
    \phi=\frac{x}{R}.
    \label{eq:conformal-coordinate-map}
\end{equation}
Since $dh=R\,dr/r$ and $dx=R\,d\phi$, substitution into
Eq.~\eqref{eq:exponential-optical-metric} gives
\begin{equation}
    d\ell_{\mathrm{opt}}^2
    =
    n_0^2
    \left(dr^2+r^2d\phi^2\right).
    \label{eq:euclidean-polar-optical-metric}
\end{equation}
Apart from the constant factor $n_0^2$, this is the Euclidean metric in polar
coordinates.  The exponential refractive medium above a flat boundary is
therefore optically equivalent to propagation in the homogeneous Euclidean exterior of
a circle.

\subsection{Boundary and ray mapping}

The physical half-plane $h\geq 0$ is mapped to
\begin{equation}
    r\geq R,
\end{equation}
while the opaque flat boundary $h=0$ becomes the circle $r=R$.  Because
multiplication of a metric by a positive constant does not alter its
geodesics, rays in the transformed coordinates are ordinary straight lines.
A limiting ray in the original medium maps to a line tangent to the circle.
The apparent horizon generated by the refractive profile is therefore not
merely a local analogy: it is the image, under Eq.~\eqref{eq:conformal-coordinate-map},
of ordinary Euclidean occlusion by a circular boundary.

Let the tangency point correspond to $(x_m,h=0)$, or equivalently to polar
angle $\phi_m=x_m/R$.  A line tangent to $r=R$ at that point satisfies
\begin{equation}
    r\cos\!\left(\phi-\phi_m\right)=R.
    \label{eq:tangent-line-polar-form}
\end{equation}
Using Eq.~\eqref{eq:conformal-coordinate-map}, this becomes
\begin{equation}
    e^{h/R}
    \cos\!\left(\frac{x-x_m}{R}\right)=1,
\end{equation}
and hence
\begin{equation}
    h(x)
    =
    -R\ln\!\left[
        \cos\!\left(\frac{x-x_m}{R}\right)
    \right].
    \label{eq:conformal-grazing-ray}
\end{equation}
Equation~\eqref{eq:conformal-grazing-ray} is precisely the grazing member of
the exact ray family in Eq.~\eqref{eq:exact-exponential-ray}.

For a point at height $h$, the angular displacement between the point and the
tangency radius follows directly from the right triangle formed by the circle
centre, the tangency point, and the endpoint:
\begin{equation}
    \Delta\phi
    =
    \arccos\!\left(\frac{R}{r}\right)
    =
    \arccos\!\left(e^{-h/R}\right).
    \label{eq:conformal-horizon-angle}
\end{equation}
Since $x=R\phi$, the corresponding physical horizontal distance is
\begin{equation}
    d(h)
    =
    R\arccos\!\left(e^{-h/R}\right),
\end{equation}
which recovers Eq.~\eqref{eq:exact-exponential-horizon-distance} for
$\kappa=1/R$ without further integration.

\subsection{Meaning and limits of the equivalence}

The equivalence is exact for the optical metric and its ray paths.
It explains, in a single construction, why the exponential profile produces
bottom-first disappearance, a limiting tangent ray, and recovery of hidden
parts when the observer or target is raised.  It also identifies these effects
with boundary occlusion in the transformed
geometry, rather than with loss of angular resolution or a generic consequence
of perspective.

The coordinate map does not, however, equate the physical height $h$ above
the flat boundary with the ordinary radial height above the circle.  The
latter is
\begin{equation}
    r-R
    =
    R\left(e^{h/R}-1\right)
    =
    h+\frac{h^2}{2R}+O\!\left(\frac{h^3}{R^2}\right).
    \label{eq:radial-height-map}
\end{equation}
Consequently, the refractive model is exactly equivalent to circular
geometry after a nonlinear redefinition of height, while comparison at the
same numerical physical height is only asymptotic for $h/R\ll 1$.  The next
section quantifies this distinction by comparing the exact refractive and
spherical horizon laws.
